\documentclass[11pt]{article}

\usepackage[margin=1in]{geometry}
\usepackage{amsmath,amssymb,mathtools}
\usepackage{booktabs,longtable,array}
\usepackage{enumitem}
\usepackage{microtype}
\usepackage[hidelinks]{hyperref}
\usepackage{xcolor}

\newcommand{\mdot}{\dot m}
\newcommand{\Tsa}{T_{\mathrm{sa}}}
\newcommand{\Tz}{T_z}
\newcommand{\QX}{Q_{\mathrm{HVAC},X}}
\newcommand{\QAC}{Q_{\mathrm{AC}}}
\newcommand{\PHVAC}{P_{\mathrm{HVAC}}}

\title{PINODE/EPSR MPC--EnergyPlus Actuator Realization\\
Simplified Frozen Paper Contract}
\author{ScaleBridge / PINODE--EPSR}
\date{29 August 2026 --- Frozen Actuator Contract v1}

\begin{document}
\maketitle

\begin{abstract}
This document freezes the deliberately simple actuator-realization interface
used to connect the PINODE/EPSR model-predictive controller (MPC) to the real
EnergyPlus 24.1 RestaurantFastFood HVAC plant.  The paper objective is not
perfect low-level actuator inversion.  It is to show, reproducibly and with
quantified evidence, that MPC commands expressed in the physically meaningful
coordinates supply-air mass flow $\mdot^\star$ and supply-air temperature
$\Tsa^\star$ can drive the actual EnergyPlus packaged single-zone (PSZ)
systems closely enough for closed-loop building control.  Directed experiments
C1--C2.7 establish an exact plant-specific fan-flow inverse, a simple
physics-based unitary sensible-load command, an empirically validated moderate
operating envelope, and one explicitly excluded extreme-cold Kitchen-heating
regime.  The final changing-command experiment passed all 36 controlled
300-s intervals.
\end{abstract}

\section{Scope and paper claim}

The final actuator layer is intentionally not a research contribution by
itself.  The actuator claim is limited to the following:

\begin{quote}
For the accepted EnergyPlus 24.1 RestaurantFastFood/Buffalo plant and the
experimentally validated operating region, physical MPC requests
$(\mdot^\star,\Tsa^\star)$ can be translated into native EnergyPlus fan-flow
and unitary sensible-load commands using a simple deterministic adapter.  The
realized 300-s HVAC response tracks the requested physical quantities with
sufficient accuracy for the paper's learned-model closed-loop MPC experiments.
\end{quote}

The paper does \emph{not} claim:
\begin{itemize}[leftmargin=2em]
\item a universal EnergyPlus actuator inverse;
\item perfect instantaneous $\Tsa$ tracking at every internal HVAC iteration;
\item exact realizability of arbitrary $(\mdot,\Tsa)$ pairs;
\item resolution of all native PSZ cycling modes;
\item a learned neural actuator model;
\item a hybrid low-level MPC for equipment cycling.
\end{itemize}

\section{Authoritative EnergyPlus plant}

The accepted EnergyPlus 24.1 candidate is
\begin{verbatim}
RestaurantFastFood_Buffalo_EPlus24.1_FINAL_CANDIDATE.idf
\end{verbatim}
with
\begin{verbatim}
SHA256 =
da86d8c3c782424d2d4f30e29a97220c9b9042f721e83591f43762f34df77101
\end{verbatim}

The paper plant contains two independently served physical zones:
\begin{equation}
\mathcal Z=\{D,K\},
\end{equation}
where $D$ is Dining and $K$ is Kitchen.

The relevant equipment is:
\begin{verbatim}
Dining:
  AirLoopHVAC:UnitarySystem  PSZ-AC_1:1
  Fan:SystemModel            PSZ-AC_1:1_addAQ Fan
  Heating coil               PSZ-AC_1:1_HeatC
  DX cooling coil            PSZ-AC_1:1_CoolC DXCoil
  Direct-air / zone inlet    PSZ-AC_1:1 Zone Equipment Inlet Node

Kitchen:
  AirLoopHVAC:UnitarySystem  PSZ-AC_2:2
  Fan:SystemModel            PSZ-AC_2:2_addAQ Fan
  Heating coil               PSZ-AC_2:2_HeatC
  DX cooling coil            PSZ-AC_2:2_CoolC DXCoil
  Direct-air / zone inlet    PSZ-AC_2:2 Zone Equipment Inlet Node
\end{verbatim}

The unitary systems are native \texttt{Load}-controlled systems.  Outdoor-air,
coil, fan, and unitary-controller physics remain native EnergyPlus physics.

\section{Closed-loop control timescale}

The MPC control interval is
\begin{equation}
\boxed{\Delta t_{\mathrm{MPC}}=300~\mathrm{s}.}
\end{equation}

EnergyPlus may execute multiple HVAC system timesteps or internal controller
iterations inside one 300-s control interval.  Therefore the paper distinguishes
instantaneous system-timestep quantities from the 300-s equivalent delivered
thermal response.

\section{Physical MPC control variables}

For each physical zone $i\in\mathcal Z$, the upper-level MPC produces
\begin{equation}
\boxed{
u_{i,k}^{\star}
=
\begin{bmatrix}
\mdot_{i,k}^{\star}\\
T_{\mathrm{sa},i,k}^{\star}
\end{bmatrix}.
}
\label{eq:physical-control}
\end{equation}

The MPC therefore remains in physically interpretable coordinates.  It does
not optimize EnergyPlus API handles or native low-level actuator coordinates.

The desired raw air-side sensible heat is
\begin{equation}
\boxed{
Q_{X,i,k}^{\star}
=
\mdot_{i,k}^{\star} c_p
\left(
T_{\mathrm{sa},i,k}^{\star}
-
T_{z,i,k}
\right).
}
\label{eq:qstar}
\end{equation}

The final implementation shall use one frozen numerical $c_p$ value
consistently in the actuator adapter, MPC, thermal-input conversion, and
reported metrics.  The screening experiments used
$c_p=1006~\mathrm{J/(kg\,K)}$.

\section{Frozen simple EnergyPlus adapter}

\subsection{Fan-flow command}

For the present two-mode PSZ/Fan:SystemModel configuration, experiment C2.2
established
\begin{equation}
\boxed{
\mdot_i^{\mathrm{EP}} = 2 a_{m,i}
}
\end{equation}
through the tested design-flow range and across Dining/Kitchen heating and
cooling.

Hence the frozen fan inverse is
\begin{equation}
\boxed{
a_{m,i,k}
=
\frac{1}{2}\mdot_{i,k}^{\star}.
}
\label{eq:fan-inverse}
\end{equation}

This mapping remained exact under simultaneous thermal actuation and under the
final changing-command trajectory test.

\subsection{Thermal command}

For the validated regimes, no learned or calibrated actuator map is retained.
The unitary sensible-load command is simply
\begin{equation}
\boxed{
a_{q,i,k}
=
Q_{X,i,k}^{\star}
=
\mdot_{i,k}^{\star}c_p
\left(
T_{\mathrm{sa},i,k}^{\star}
-
T_{z,i,k}
\right).
}
\label{eq:simple-thermal-command}
\end{equation}

Thus the final paper actuator adapter is only
\begin{equation}
\boxed{
\begin{aligned}
a_{m,i,k}
&=
\tfrac12\mdot_{i,k}^{\star},\\
a_{q,i,k}
&=
\mdot_{i,k}^{\star}c_p
\left(
T_{\mathrm{sa},i,k}^{\star}
-
T_{z,i,k}
\right).
\end{aligned}}
\label{eq:frozen-adapter}
\end{equation}

\subsection{EnergyPlus callback chronology}

The final callback contract is
\begin{verbatim}
Fan Air Mass Flow Rate
    -> InsideHVACSystemIterationLoop

Unitary HVAC / Sensible Load Request
    -> AfterPredictorAfterHVACManagers
\end{verbatim}

C2.5 found numerically identical physical responses when the sensible-load
request was instead issued from the HVAC-system iteration callback.  Therefore
the simpler \texttt{AfterPredictorAfterHVACManagers} chronology is retained.

\section{Text schematic of the frozen architecture}

\begin{verbatim}
weather / disturbances / price / grid reference
                    |
                    v
        +-----------------------+
        | learned thermal model |
        | RC / PINN-RC / NODE   |
        | Base PINODE / EBP     |
        +-----------+-----------+
                    |
                    v
                 MPC
                    |
                    | physical request
                    | [m_dot*, T_sa*]
                    v
        +-----------------------+
        | SIMPLE FROZEN ADAPTER |
        |                       |
        | a_fan = 0.5 m_dot*    |
        |                       |
        | Q_X* = m_dot* cp      |
        |       (T_sa* - T_z)   |
        |                       |
        | a_q = Q_X*            |
        +-----------+-----------+
                    |
                    v
        +-----------------------+
        | EnergyPlus 24.1 PSZ   |
        | native OA / coils /   |
        | fan / unitary physics |
        +-----------+-----------+
                    |
                    v
          m_dot_EP , T_sa_EP
                    |
          +---------+---------+
          |                   |
          v                   v
   delivered Q_X          next T_zone
          |                   |
          v                   |
       Phase-C G_Q            |
          |                   |
          v                   |
        Q_AC -----------------+
          |
          v
       thermal ODE

Separate electrical path:
    |Q_AC| -> Phase-C G_P -> P_HVAC -> MPC energy/price/grid objective
\end{verbatim}

\section{300-s equivalent supply temperature}

The final actuator acceptance is evaluated at the same 300-s timescale used by
the MPC and thermal models.

For one control interval $k$, define the time-weighted delivered raw sensible
heat
\begin{equation}
\overline{Q}_{X,i,k}^{\mathrm{EP}}
=
\frac{
\sum_r Q_{X,i,r}^{\mathrm{EP}}\Delta t_r
}{
\sum_r \Delta t_r
},
\end{equation}
with corresponding time-weighted zone temperature and mass flow
$\overline{T}_{z,i,k}$ and $\overline{\mdot}_{i,k}^{\mathrm{EP}}$.

The equivalent supply temperature is
\begin{equation}
\boxed{
T_{\mathrm{sa,eq},i,k}^{\mathrm{EP}}
=
\overline{T}_{z,i,k}
+
\frac{
\overline{Q}_{X,i,k}^{\mathrm{EP}}
}{
\overline{\mdot}_{i,k}^{\mathrm{EP}} c_p
}.
}
\label{eq:tsa-eq}
\end{equation}

This does not hide the native system-timestep response.  The closed-loop
experiment should log both the detailed EnergyPlus system-timestep quantities
and the 300-s equivalent quantities.

\section{Empirically supported operating envelope}

The final trajectory smoke directly exercised moderate commands with
\begin{equation}
\boxed{
0.50
\le
\frac{\mdot_i^\star}{\mdot_{i,\max}}
\le
0.80
}
\label{eq:flow-envelope}
\end{equation}
and
\begin{equation}
\boxed{
4^\circ\mathrm C
\le
\left|
T_{\mathrm{sa},i}^\star - T_{z,i}
\right|
\le
8^\circ\mathrm C
}
\label{eq:thermal-envelope}
\end{equation}
for the validated heating/cooling trajectories.

This is an empirically validated paper operating envelope, not an exhaustive
proof that every point in the Cartesian rectangle is realizable under every
weather and native-controller state.

Define the paper actuator-feasible set
\begin{equation}
\mathcal U_i^{\mathrm{act}}
(\xi_i)
\end{equation}
as the subset of physical commands supported by this evidence under local HVAC
context $\xi_i$.  The final MPC command set must satisfy
\begin{equation}
\boxed{
u_{i,k}^{\star}
\in
\mathcal U_i^{\mathrm{MPC}}
=
\mathcal U_i^{\mathrm{train}}
\cap
\mathcal U_i^{\mathrm{act}}.
}
\label{eq:final-feasible-set}
\end{equation}

The intersection with the thermal-model TRAIN support is important: the paper
should not ask the learned model to predict far outside the control/input
region represented in its training data.

\section{Generic symmetric feasibility supervisor and Kitchen-heating evidence}

All four physical zone/mode combinations are commandable:
\begin{center}
\begin{tabular}{lcc}
\toprule
Zone & Heating & Cooling\\
\midrule
Dining & MPC-commandable & MPC-commandable\\
Kitchen & MPC-commandable & MPC-commandable\\
\bottomrule
\end{tabular}
\end{center}

There is no categorical Kitchen-heating exclusion in the final simulator.
Instead, Dining and Kitchen use the same numerical per-zone supervisor and the
same actuator transformation.

At each 300-s control boundary, define
\begin{align}
\gamma_{\dot m,i}
&=
\frac{\mdot_i^\star}{\mdot_{i,\max}},
\\
\Delta T_i^\star
&=
T_{\mathrm{sa},i}^\star-T_{z,i}.
\end{align}
The current configured supervisor accepts an external override when
\begin{equation}
\boxed{
0.50\le\gamma_{\dot m,i}\le0.80,
\qquad
4^\circ\mathrm C
\le
|\Delta T_i^\star|
\le
8^\circ\mathrm C.
}
\label{eq:supervisor-envelope}
\end{equation}
The sign of \(\Delta T_i^\star\) is unrestricted.

If zone \(i\) is outside the configured envelope, the simulator releases both
the fan-flow and unitary sensible-load overrides for that zone for the current
300-s interval:
\begin{equation}
\boxed{
u_{i,k}^{\star}\notin\mathcal U_i^{\mathrm{sup}}
\Longrightarrow
\text{native EnergyPlus HVAC for zone }i.
}
\label{eq:per-zone-fallback}
\end{equation}
The other zone remains under external control when its own command is accepted.

The extreme-cold Kitchen-heating behavior observed in C2.3--C2.6 remains
scientifically important.  Independent forced airflow combined with native
fuel-coil ON/OFF behavior produced alternating hot and very cold supply-air
substeps and poor 300-s \(T_{\mathrm{sa}}\) tracking in those particular
conditions.  C2.6 also showed that an exploratory affine correction did not
repair this behavior.  Those results justify retaining detailed
commanded-versus-realized history and a native-fallback mechanism; they do
\textbf{not} justify a blanket prohibition on Kitchen heating.

The final controller therefore does not add a learned actuator model or hybrid
duty-cycle inverse.  All learned-model families use the same simple adapter,
the same supervisor, and the same per-zone fallback mechanism.

\section{Evidence ladder C1--C2.7}

\begin{longtable}{p{0.10\textwidth}p{0.25\textwidth}p{0.55\textwidth}}
\toprule
Stage & Purpose & Frozen conclusion \\
\midrule
\endfirsthead
\toprule
Stage & Purpose & Frozen conclusion \\
\midrule
\endhead

C1 &
Actuator/API inventory &
Established the EnergyPlus actuator/variable registry needed for directed
experiments.  Broad actuator discovery was then closed. \\

C2.1 &
Directed causal screening &
Terminal mass-flow and generic supply-node temperature-setpoint paths were
rejected for this topology.  Fan mass-flow, unitary sensible-load request, and
DX-speed paths showed causal effects. \\

C2.2 &
Fan calibration &
For Dining/Kitchen and heating/cooling,
$\mdot^{EP}=2a_m$ with ratio error zero up to numerical roundoff across tested
flow levels.  Freeze $a_m=\tfrac12\mdot^\star$. \\

C2.3 &
Fixed-flow thermal screen &
With airflow fixed, direct physics
$a_q=\mdot^\star c_p(\Tsa^\star-\Tz)$ was already promising for Dining heating,
Dining cooling, and Kitchen cooling.  Kitchen extreme-cold heating exposed
native cycling/makeup-air complications. \\

C2.4 &
Counterfactual characterization &
Independent counterfactual restarts showed strong local active-regime
relationships while also revealing that interval averages can be dominated by
native ON/OFF unitary cycling.  No complex map was frozen. \\

C2.5 &
Callback chronology closure &
\texttt{AfterPredictorAfterHVACManagers} and
\texttt{InsideHVACSystemIterationLoop} produced numerically identical physical
responses for the tested sensible-load cases.  Freeze the simpler
AfterPredictor chronology. \\

C2.6 &
Held-out static validation &
Direct physics passed six of eight held-out cases.  Kitchen heating at 50\%
and 80\% design flow exposed difficult native cycling and poor realized
supply-temperature tracking.  The exploratory affine correction failed and
was rejected; the evidence is retained as a robustness warning rather than a
categorical mode ban. \\

C2.7 &
Final changing-trajectory smoke &
With only the final simple adapter, simultaneous Dining+Kitchen cooling and
Dining heating passed all 36 controlled 300-s intervals.  Static actuator
development is frozen. \\

Simulator closure &
Direct Runtime-API integration &
The project-owned EnergyPlus simulator exercises full 300-s Dining/Kitchen
cooling and heating command transport, symmetric per-zone feasibility
decisions, native fallback, and broad history capture.  This runtime closure
does not replace the C2.6/C2.7 physical-tracking evidence. \\
\bottomrule
\end{longtable}

\section{Key numerical evidence}

\subsection{C2.2 fan-flow inverse}

For all four zone/mode combinations, the median realized-to-command flow ratio
was
\begin{equation}
\boxed{2.0}
\end{equation}
with maximum absolute deviation from the factor-two relationship equal to zero
or floating-point roundoff ($4.44\times10^{-16}$).

\subsection{C2.5 moderate physical tracking}

At $\mdot^\star=0.66\mdot_{\max}$ and
$|T_{\mathrm{sa}}^\star-T_z|=8^\circ$C:
\begin{center}
\begin{tabular}{llrr}
\toprule
Zone & Mode & Realized $\Delta T$ [$^\circ$C] & $T_{\mathrm{sa}}$ MAE [$^\circ$C] \\
\midrule
Dining & heating & +7.789 & 0.211 \\
Dining & cooling & -7.859 & 0.141 \\
Kitchen & cooling & -8.132 & 0.141 \\
Kitchen & heating & -5.334 & 13.373 \\
\bottomrule
\end{tabular}
\end{center}
Flow MAE and sensible-load actuator-readback MAE were zero in all four cases.

\subsection{C2.6 held-out static validation}

C2.6 moved away from the thermal-identification point by testing 50\% and 80\%
design flow with a held-out $6^\circ$C target:
\begin{center}
\begin{tabular}{llrrl}
\toprule
Zone & Mode & Flow & $T_{\mathrm{sa}}$ MAE [$^\circ$C] & Gate \\
\midrule
Dining & heating & 0.50 & 0.291 & PASS \\
Dining & heating & 0.80 & 0.138 & PASS \\
Dining & cooling & 0.50 & 0.243 & PASS \\
Dining & cooling & 0.80 & 0.022 & PASS \\
Kitchen & cooling & 0.50 & 0.319 & PASS \\
Kitchen & cooling & 0.80 & 0.268 & PASS \\
Kitchen & heating & 0.50 & 25.957 & FAIL \\
Kitchen & heating & 0.80 & 23.709 & FAIL \\
\bottomrule
\end{tabular}
\end{center}
All eight held-out cases retained exact airflow tracking.

\subsection{C2.7 final actuator acceptance}

The final changing-command experiment used only
Eqs.~\eqref{eq:fan-inverse}--\eqref{eq:simple-thermal-command}.  It contained
simultaneous two-zone cooling and a separate Dining-heating trajectory.

The scientific acceptance result was:
\begin{align}
N_{\mathrm{controlled}}
&=36,\\
N_{\mathrm{pass}}
&=36,\\
\text{pass fraction}
&=1.000,\\
\operatorname{mean}
\left|
T_{\mathrm{sa,eq}}^{EP}
-
T_{\mathrm{sa}}^\star
\right|
&=
0.378955^\circ\mathrm C,\\
\max
\left|
T_{\mathrm{sa,eq}}^{EP}
-
T_{\mathrm{sa}}^\star
\right|
&=
1.626835^\circ\mathrm C,\\
\max
\left|
\frac{
\mdot^{EP}-\mdot^\star
}{
\mdot^\star
}
\right|
&=
0,\\
\text{all requested thermal signs correct}
&=
\mathrm{true}.
\end{align}

The predeclared gates were:
\begin{align}
\text{minimum interval pass fraction} & = 0.90,\\
\text{maximum mean }|T_{\mathrm{sa,eq}}\text{ error}| & = 1.5^\circ\mathrm C,\\
\text{per-interval }|T_{\mathrm{sa,eq}}\text{ error}| & \le 2^\circ\mathrm C,\\
\text{flow gate} & = 0.5\%\text{ of design flow}.
\end{align}

Therefore
\begin{equation}
\boxed{\text{FINAL ACTUATOR ACCEPT}=\mathrm{TRUE}.}
\end{equation}

\section{Known EnergyPlus baseline diagnostics}

The actuator experiments repeatedly observed two
\texttt{CheckAirLoopFlowBalance} severe markers associated with
PSZ-AC\_1:1 and PSZ-AC\_2:2.  These were previously investigated during the
EnergyPlus 9.0$\rightarrow$24.1 transition and accepted as documented
version-dependent/transient airflow-controller diagnostics for the unchanged
candidate.  C2.5/C2.7 introduced no new Fatal actuator-related failures.

The original EnergyPlus 9.0 model and accepted EnergyPlus 24.1 candidate must
remain unchanged by the actuator adapter.  All experiments are scratch-only.

\section{Final direct EnergyPlus closed-loop interface}

The final software stack uses a project-owned synchronous wrapper around the
EnergyPlus~24.1 Python Runtime/Data Exchange API.  Sinergym and Gymnasium are
not required by the paper closed-loop path.  The wrapper owns callback
synchronization, feasibility supervision, actuator realization, and broad
history capture.

\subsection*{Interface to the learned thermal models}

When model training is complete, the closed-loop experiment shall use the same
frozen actuator adapter for every thermal-model family.

At MPC time $k$:
\begin{enumerate}[leftmargin=2em]
\item estimate/initialize the model state $\hat x_k$ using the frozen
method-specific observer/encoder;
\item solve the upper-level MPC for physical
$u_{i,k}^\star=[\mdot_{i,k}^\star,T_{\mathrm{sa},i,k}^\star]^\top$;
\item enforce the final model-support bounds in the MPC and evaluate the
generic per-zone simulator supervisor of Eq.~\eqref{eq:supervisor-envelope};
\item for an accepted zone command, compute $a_{m,i,k}$ and $a_{q,i,k}$ from
Eq.~\eqref{eq:frozen-adapter}; if rejected, release both overrides for that
zone and use native EnergyPlus HVAC for the interval;
\item issue the two commands at their frozen EnergyPlus callbacks;
\item advance EnergyPlus for 300 s;
\item record both detailed system-timestep and 300-s equivalent actuator
response;
\item update measurements/state estimate and repeat.
\end{enumerate}

The paper comparison must hold constant across model families:
\begin{itemize}[leftmargin=2em]
\item EnergyPlus candidate and EPW;
\item initial date/time and plant state;
\item forecasts and disturbances;
\item price/grid signal;
\item comfort constraints;
\item MPC horizon and solver settings except method-specific necessities;
\item physical input bounds and actuator-feasible region;
\item actuator adapter and fallback;
\item Phase-C $G_Q$ and $G_P$ mappings;
\item evaluation metrics.
\end{itemize}

\section{Return point after learned-model development}

Actuator development is now frozen.  The immediate paper-development sequence
is:
\begin{verbatim}
1. Train / validate all learned thermal-model methods.
2. Freeze comparable model checkpoints and observers.
3. Return to closed-loop MPC.
4. Use THIS actuator contract unchanged.
5. Run common EnergyPlus closed-loop comparison.
6. Report comfort, energy/cost/grid, model, optimization,
   and commanded-vs-realized actuator metrics.
\end{verbatim}

No additional actuator identification should be performed unless the actual
closed-loop experiment exposes a concrete failure inside the already validated
operating envelope.

\section{Reproducibility evidence archive}

The authoritative LaTeX contracts are repository-owned and live under
\texttt{Paper\_PINODE\_EPSR/docs/mathematics/contracts}.  Scientific
EnergyPlus evidence remains under the separate EPSR data root.

The frozen scientific evidence archive should retain:
\begin{enumerate}[leftmargin=2em]
\item a manifest identifying the repository-owned mathematical contract version;
\item all C1--C2.7 reproduction scripts;
\item all portable PASS evidence ZIPs available from the experiments;
\item snapshots of the persistent
\texttt{03\_mpc\_actuator\_audit/Day7C*} directories, including failed or
superseded attempts where present;
\item SHA256 manifests for scripts, portable bundles, and local evidence
snapshots;
\item the historical detailed Part-7 exploratory mathematical contract.
\end{enumerate}

The exploratory affine and hybrid interpretations are retained as scientific
history but are not part of the frozen controller.

\section{Final frozen statement}

For the RestaurantFastFood identity two-zone plant, the physical MPC command is
\begin{equation}
\boxed{
u_k^\star
=
\begin{bmatrix}
\mdot_{D,k}^{\star}\\
T_{\mathrm{sa},D,k}^{\star}\\
\mdot_{K,k}^{\star}\\
T_{\mathrm{sa},K,k}^{\star}
\end{bmatrix}.
}
\end{equation}

For either zone \(i\in\{D,K\}\), an accepted command is realized through
\begin{equation}
\boxed{
\begin{aligned}
a_{m,i,k}
&=
\tfrac12\mdot_{i,k}^{\star},\\
a_{q,i,k}
&=
\mdot_{i,k}^{\star}c_p
\left(
T_{\mathrm{sa},i,k}^{\star}
-
T_{z,i,k}
\right),
\qquad c_p=1006~\mathrm{J/(kg\,K)}.
\end{aligned}}
\end{equation}

Dining heating, Dining cooling, Kitchen heating, and Kitchen cooling are all
MPC-commandable.  The simulator applies the same numerical feasibility
supervisor to every zone/mode.  If one zone is outside the configured
supervisor envelope, both external overrides for that zone are released and
native EnergyPlus HVAC controls that zone for the current 300-s interval; the
other zone is unaffected.

The C2.7 physical-trajectory test remains the actuator-tracking acceptance
evidence: 36/36 changing 300-s controlled intervals passed, with zero measured
flow error, mean equivalent supply-temperature error
$0.379^\circ$C, maximum error $1.627^\circ$C, and correct requested thermal
sign in every controlled interval.  C2.6 Kitchen-heating failures remain
reported as regime-specific physical evidence and motivate explicit
commanded-versus-realized logging and fallback, not a categorical mode ban.

The final software plant boundary is the direct EnergyPlus~24.1 Runtime/Data
Exchange API through the project-owned generic simulator.  Sinergym is not a
required dependency.

\end{document}
